% SOURCE: Chinese PDF viewer pp. 72--76.
\section{Ensembles and the Density Operator}
We consider the spin properties of a many-particle system.

\subsection*{1) Pure ensembles and mixed ensembles}
Identical physical systems form an ensemble; for example, an ensemble may consist of particles with identical spins. The spin state of a spin-$1/2$ particle, with azimuthal and polar angles $\alpha,\beta$, is
\[\chi(\alpha,\beta)=c_+(\alpha,\beta)\chi_++c_-(\alpha,\beta)\chi_-
=\cos\frac\beta2e^{-\ii\alpha/2}\chi_+
+\sin\frac\beta2e^{\ii\alpha/2}\chi_-,\]
where $\chi_+,\chi_-$ are eigenstates of $\op s_z$ and
\[\frac{c_+}{c_-}=\frac{\cos(\beta/2)}{e^{\ii\alpha}\sin(\beta/2)}.\]
If all particles are in the same state, e.g., their spins point in the same direction, the ensemble is called a \emph{pure ensemble} or a polarized system. If particles are in different states, e.g., their spins do not all point in one direction, the ensemble is called a \emph{mixed ensemble}. For example, if $70\%$ of particles have spin up and $30\%$ spin down, the system is partially polarized and forms a partially mixed ensemble. In a system with completely random spin directions, the probabilities of spin up and spin down are each $50\%$. Their overall contributions cancel and there is no net spin polarization.

\subsection*{2) Ensemble averages and the density operator}
For a fixed state $\ket\alpha$, the average of a mechanical quantity $\op A$ is
\[\langle A\rangle_\alpha=\mel\alpha{\op A}\alpha.\]
This is a \emph{quantum average}. Insert any complete basis $B$:
\[\langle A\rangle_\alpha
=\sum_{b,b'}\braket\alpha{b'}\mel{b'}{\op A}b\braket b\alpha.\]
The ensemble average is
\[[A]=\sum_\alpha w_\alpha\langle A\rangle_\alpha,\]
where $w_\alpha$ is the probability that the system occupies state $\ket\alpha$. It satisfies the normalization condition
\[\sum_\alpha w_\alpha=1.\]
This is a \emph{statistical average}. Since
\begin{align*}
[A]&=\sum_{\alpha,b,b'}w_\alpha\braket\alpha{b'}\mel{b'}{\op A}b\braket b\alpha\\
&=\sum_{b,b'}\left(\sum_\alpha w_\alpha\braket b\alpha\braket\alpha{b'}\right)
\mel{b'}{\op A}b,
\end{align*}
we define the density operator by
\[\boxed{\op\rho=\sum_\alpha w_\alpha\ket\alpha\bra\alpha.}\]
Its matrix elements in the $B$ representation are
\[\rho_{bb'}=\mel b{\op\rho}{b'}
=\sum_\alpha w_\alpha\braket b\alpha\braket\alpha{b'}.\]
Therefore
\[\boxed{[A]=\sum_{b,b'}\mel b{\op\rho}{b'}\mel{b'}{\op A}b
=\sum_b\mel b{\op\rho\op A}b
\equiv\operatorname{tr}(\op\rho\op A).}\]
This fundamental equation of quantum statistical mechanics includes both quantum and statistical averaging. The trace is unchanged by a change of representation, so the formula is independent of the chosen representation.

In a continuous representation, e.g., the position basis, the matrix elements are
\[\rho_{\vect x\vect x'}=\mel{\vect x}{\op\rho}{\vect x'}
=\sum_\alpha w_\alpha\braket{\vect x}\alpha\braket\alpha{\vect x'}
=\sum_\alpha w_\alpha\psi_\alpha(\vect x)\psi_\alpha^*(\vect x').\]
The ensemble average becomes
\[[A]=\operatorname{tr}(\op\rho\op A)
=\int\dd^3x\,\mel{\vect x}{\op\rho\op A}{\vect x}.\]

The density matrix satisfies the normalization condition
\begin{align*}
\operatorname{tr}\op\rho
&=\sum_{\alpha,b}w_\alpha\braket b\alpha\braket\alpha b\\
&=\sum_\alpha w_\alpha\bra\alpha\left(\sum_b\ket b\bra b\right)\ket\alpha
&&\text{(completeness)}\\
&=\sum_\alpha w_\alpha\braket\alpha\alpha
&&\text{(quantum-state normalization)}\\
&=\sum_\alpha w_\alpha=1
&&\text{(statistical normalization)}.
\end{align*}
Suppose the density operator has an eigenstate $\ket\theta$:
\[\op\rho\ket\theta=\theta\ket\theta,\qquad
\op\rho^2\ket\theta=\theta^2\ket\theta.\]
For a pure ensemble all systems are in the same state $\ket n$:
\[w_\alpha=\begin{cases}1,&\alpha=n,\\0,&\alpha\ne n.\end{cases}\]
Then
\[\op\rho=\ket n\bra n,\qquad
\op\rho^2=\ket n\braket n n\bra n=\ket n\bra n=\op\rho,\]
so $\theta^2=\theta$ and its eigenvalues are $\theta=0,1$. If $\ket n$ is a basis state of the chosen representation, the density matrix is diagonal:
\[\rho_{bb'}=\braket b n\braket n{b'}=\delta_{bn}\delta_{b'n},\qquad
\rho=\begin{pmatrix}0&&&&\\&\ddots&&&\\&&1&&\\&&&\ddots&\\&&&&0\end{pmatrix}.\]
Only one matrix entry is $1$; all others are $0$.

For a completely mixed ensemble, all possible states occur with equal probabilities $w_\alpha=w_0=1/N$, where $N$ is the number of states. Then
\[\rho_{bb'}=\sum_\alpha w_\alpha\braket b\alpha\braket\alpha{b'}
=\frac1N\braket b{b'}=\frac1N\delta_{bb'},\]
and the density matrix is proportional to the identity:
\[\rho=\frac1N\begin{pmatrix}1&&\\&\ddots&\\&&1\end{pmatrix}.\]

\textbf{Example 1: pure ensemble.} Suppose that the completely polarized state is $\ket n=\ket+$, with $\op s_z\ket\pm=\pm(\hbar/2)\ket\pm$. Then
\[\op\rho=\ket+\bra+.\]
In the $s_z$ basis $\rho$ is a $2\times2$ matrix:
\[\rho_{++}=\braket+{+}\braket+{+}=1,\quad
\rho_{+-}=\braket+{+}\braket+{-}=0,\]
\[\rho_{-+}=\braket-{+}\braket+{+}=0,\quad
\rho_{--}=\braket-{+}\braket+{-}=0,\qquad
\rho=\begin{pmatrix}1&0\\0&0\end{pmatrix}.\]
If the completely polarized state is
\[\ket n=\ket{s_x^+}=\frac1{\sqrt2}\ket++\frac1{\sqrt2}\ket-,\qquad
\op s_x\ket{s_x^+}=\frac\hbar2\ket{s_x^+},\]
then $\op\rho=\ket{s_x^+}\bra{s_x^+}$. In the $s_z$ representation,
\[\rho_{++}=\braket+{s_x^+}\braket{s_x^+}{+}=\frac12,\quad
\rho_{+-}=\braket+{s_x^+}\braket{s_x^+}{-}=\frac12,\]
\[\rho_{-+}=\braket-{s_x^+}\braket{s_x^+}{+}=\frac12,\quad
\rho_{--}=\braket-{s_x^+}\braket{s_x^+}{-}=\frac12,\qquad
\rho=\frac12\begin{pmatrix}1&1\\1&1\end{pmatrix}.\]

\textbf{Example 2: density matrix for a completely unpolarized system.} With $N=2$,
\[\op\rho=\frac12\ket+\bra++\frac12\ket-\bra-.\]
In the $s_z$ representation,
\[\rho_{++}=\frac12,\quad\rho_{+-}=0,\quad\rho_{-+}=0,\quad\rho_{--}=\frac12,\qquad
\rho=\frac12\begin{pmatrix}1&0\\0&1\end{pmatrix}.\]
The density matrix of a completely mixed ensemble is proportional to the identity. The expectation values in a completely unpolarized ensemble obviously cancel:
\[[s_i]=\operatorname{tr}(\op\rho\op s_i)=\frac12\operatorname{tr}(\op s_i).\]
Inserting the spin matrices $s_x,s_y,s_z$ gives $[s_i]=0$.

\textbf{Example 3: density matrix for a partially polarized system.} Let
\[\op\rho=\frac34\ket+\bra++\frac14\ket{s_x^+}\bra{s_x^+}.\]
In the $s_z$ representation,
\[\rho_{++}=\frac34+\frac14\frac12=\frac78,\qquad
\rho_{+-}=\rho_{-+}=\rho_{--}=\frac18,\]
so that
\[\rho=\frac18\begin{pmatrix}7&1\\1&1\end{pmatrix},\qquad
[s_x]=\frac\hbar8,\quad[s_y]=0,\quad[s_z]=\frac{3\hbar}{8}.\]

\subsection*{3) Evolution of the ensemble}
How does the density operator evolve in time? If the statistical properties of the ensemble do not vary in time, its probability distribution $w_\alpha$ is constant and
\[\op\rho(t)=\sum_\alpha w_\alpha\ket{\alpha,t}\bra{\alpha,t}.\]
By the state evolution equations
\[\ii\hbar\frac\partial{\partial t}\ket{\alpha,t}=\op H\ket{\alpha,t},\qquad
-\ii\hbar\frac\partial{\partial t}\bra{\alpha,t}=\bra{\alpha,t}\op H,\]
we obtain
\begin{align*}
\ii\hbar\frac{\partial\op\rho}{\partial t}
&=\sum_\alpha w_\alpha\left(\op H\ket{\alpha,t}\bra{\alpha,t}
-\ket{\alpha,t}\bra{\alpha,t}\op H\right)\\
&=-[\op\rho,\op H].
\end{align*}
This is the density operator's time-evolution equation. It resembles the equation of motion of mechanical quantities in the Heisenberg picture, differing by a minus sign; however $\op\rho$ is not an observable operator but an operator constructed from states in the Schrödinger picture.